\chapter*{Conclusion générale}
\addcontentsline{toc}{chapter}{Conclusion générale}

Le projet E-biessa aboutit à une plateforme complète de mise en relation autour des pièces automobiles d'occasion. Le socle transactionnel prend en charge les comptes, les rôles, les profils, les annonces riches, les images, la compatibilité multi-véhicule, les favoris, le contact, la messagerie, la réputation et l'administration. L'expérience Angular rend ces fonctions accessibles dans une interface cohérente et responsive.


L'apport majeur est l'intégration contrôlée de l'intelligence artificielle. Le modèle adapté comprend des formulations multilingues proches du langage des utilisateurs. La recherche vectorielle améliore le rappel, tandis que la validation déterministe protège la précision. Les recommandations et les agents spécialisés utilisent les mêmes principes : lecture seule, données minimales, raisons visibles et retour systématique à la source de vérité.


Ce travail a permis de mobiliser la conception orientée domaine, la sécurité web, les architectures distribuées, l'UX, le traitement du langage et l'évaluation expérimentale. La séparation Spring Boot / FastAPI représente un compromis justifié entre robustesse métier et agilité IA.


\section{Perspectives}

\begin{itemize}
\item déployer la pile derrière HTTPS après rotation et gestion centralisée des secrets ;

\item faire valider la compatibilité par des professionnels et intégrer des catalogues OEM sous licence ;

\item constituer un benchmark plus large, annoté manuellement, couvrant dialectes et fautes d'orthographe ;

\item mesurer la qualité des recommandations avec des retours anonymisés et des métriques hors ligne ;

\item ajouter notifications, paiement ou réservation uniquement après validation juridique et opérationnelle ;

\item surveiller dérive du modèle, qualité des données et taux de signalement avant tout usage public.
\end{itemize}
